% preface.tex — Pr\'eface du cours de Th\'eorie des Points Fixes
\chapter*{Pr\'eface}
\addcontentsline{toc}{chapter}{Pr\'eface}

La th\'eorie des points fixes constitue l'un des piliers fondamentaux de l'analyse
math\'ematique moderne. Son influence s'\'etend bien au-del\`a des fronti\`eres de
l'analyse fonctionnelle, irriguant des domaines aussi vari\'es que la topologie,
l'alg\`ebre, l'optimisation, la th\'eorie des jeux, l'\'economie math\'ematique et
les \'equations aux d\'eriv\'ees partielles.

\section*{Objectifs du cours}

Ce cours a \'et\'e con\c{c}u pour les \'etudiants de Master et de Doctorat en
math\'ematiques. Il vise \`a pr\'esenter de mani\`ere rigoureuse et compl\`ete les
grands th\'eor\`emes de points fixes, leurs d\'emonstrations, leurs g\'en\'eralisations
et leurs applications.

Les objectifs p\'edagogiques principaux sont~:
\begin{itemize}
    \item Ma\^itriser les d\'emonstrations compl\`etes des th\'eor\`emes classiques
          (Banach, Brouwer, Schauder, Kakutani, Tarski--Knaster);
    \item Comprendre les liens profonds entre ces r\'esultats et les structures
          math\'ematiques sous-jacentes (m\'etriques, topologiques, ordonn\'ees);
    \item Savoir appliquer ces th\'eor\`emes \`a la r\'esolution de probl\`emes concrets~:
          \'equations diff\'erentielles, \'equations int\'egrales, \'equilibres de Nash;
    \item Explorer les extensions modernes~: espaces m\'etriques g\'en\'eralis\'es,
          points fixes al\'eatoires, correspondances multivoques.
\end{itemize}

\section*{Organisation du cours}

Le cours est structur\'e en onze chapitres organis\'es en trois parties naturelles.

\subsection*{Partie I~: Th\'eor\`emes m\'etriques (Chapitres 1--4)}

La premi\`ere partie est consacr\'ee aux r\'esultats de type m\'etrique. Apr\`es un
chapitre d'introduction qui pose les motivations historiques et les d\'efinitions
fondamentales, le Chapitre~2 pr\'esente le c\'el\`ebre Principe de Contraction de
Banach (1922) et ses g\'en\'eralisations directes. Le Chapitre~3 explore les
extensions m\'etriques dues \`a Edelstein, Kannan, \'Ciri\'c et Caristi, qui
rel\^achent de diff\'erentes mani\`eres l'hypoth\`ese de contraction. Le Chapitre~4
introduit les espaces m\'etriques g\'en\'eralis\'es ($b$-m\'etriques, $G$-m\'etriques)
et examine comment les th\'eor\`emes classiques s'y transposent.

\subsection*{Partie II~: Th\'eor\`emes topologiques (Chapitres 5--7)}

La deuxi\`eme partie aborde les th\'eor\`emes de nature topologique. Le Chapitre~5
est d\'edi\'e au th\'eor\`eme de Brouwer (1911), v\'eritable joyau de la topologie
alg\'ebrique, avec plusieurs d\'emonstrations (combinatoire via le lemme de Sperner,
analytique, homologique). Le Chapitre~6 traite du th\'eor\`eme de Schauder (1930),
g\'en\'eralisation en dimension infinie du th\'eor\`eme de Brouwer, et de ses variantes
(Schauder--Tychonoff, Darbo). Le Chapitre~7 pr\'esente le th\'eor\`eme de Kakutani
(1941) pour les correspondances multivoques, outil indispensable en th\'eorie des
jeux et en \'economie math\'ematique.

\subsection*{Partie III~: Extensions et applications (Chapitres 8--11)}

La troisi\`eme partie regroupe des th\`emes avanc\'es. Le Chapitre~8 \'etudie les
points fixes dans les espaces ordonn\'es (th\'eor\`emes de Tarski, Knaster--Tarski,
Bourbaki--Witt). Le Chapitre~9 est un chapitre transversal d'applications~:
\'equations diff\'erentielles ordinaires, \'equations aux d\'eriv\'ees partielles,
\'equilibres de Nash, mod\`eles \'economiques. Le Chapitre~10 pr\'esente le
th\'eor\`eme de Markov--Kakutani et ses cons\'equences pour les points fixes communs
de familles de transformations. Enfin, le Chapitre~11 introduit la th\'eorie des
points fixes al\'eatoires, domaine en plein essor.

\section*{Pr\'erequis}

Le lecteur doit poss\'eder des connaissances solides en~:
\begin{itemize}
    \item \textbf{Analyse fonctionnelle}~: espaces de Banach, espaces de Hilbert,
          op\'erateurs lin\'eaires born\'es, th\'eor\`eme de Hahn--Banach, compacit\'e
          faible;
    \item \textbf{Topologie g\'en\'erale}~: espaces compacts, connexit\'e,
          th\'eor\`emes de s\'eparation, espaces m\'etriques complets;
    \item \textbf{Th\'eorie de la mesure}~: espaces mesurables, fonctions mesurables,
          int\'egration de Lebesgue (pour le Chapitre~11);
    \item \textbf{Alg\`ebre lin\'eaire}~: espaces vectoriels de dimension finie,
          d\'eterminants, formes multilin\'eaires.
\end{itemize}

\section*{Notations}

Tout au long de ce cours, nous utiliserons les notations suivantes~:
\begin{itemize}
    \item $\R$, $\N$, $\Z$, $\Q$, $\C$~: les ensembles de nombres r\'eels,
          naturels, entiers, rationnels et complexes;
    \item $(X, d)$~: un espace m\'etrique, o\`u $d$ d\'esigne la distance;
    \item $\norm{\cdot}$~: une norme sur un espace vectoriel norm\'e;
    \item $\inner{\cdot}{\cdot}$~: un produit scalaire sur un espace de Hilbert;
    \item $B(x, r)$~: la boule ouverte de centre $x$ et de rayon $r$;
    \item $\overline{B}(x, r)$~: la boule ferm\'ee de centre $x$ et de rayon $r$;
    \item $\overline{A}$~: l'adh\'erence d'un ensemble $A$;
    \item $\mathrm{co}(A)$, $\overline{\mathrm{co}}(A)$~: l'enveloppe convexe et
          l'enveloppe convexe ferm\'ee de $A$;
    \item $\mathrm{Fix}(T)$~: l'ensemble des points fixes de l'application $T$;
    \item $\mathrm{Lip}(T)$~: la constante de Lipschitz de $T$.
\end{itemize}

\section*{Conventions}

Sauf mention contraire~:
\begin{itemize}
    \item Les espaces vectoriels sont sur le corps $\R$ des r\'eels;
    \item Les espaces topologiques sont suppos\'es s\'epar\'es (Hausdorff);
    \item Les applications sont \`a valeurs r\'eelles ou dans un espace de Banach;
    \item Le symbole $\square$ marque la fin d'une d\'emonstration;
    \item Le symbole $\lozenge$ marque la fin d'un exemple ou d'une remarque.
\end{itemize}

\section*{Guide de lecture}

Ce cours peut \^etre abord\'e de plusieurs mani\`eres selon le niveau et les
int\'er\^ets du lecteur~:

\begin{itemize}
    \item \textbf{Parcours fondamental} (M1)~: Chapitres 1, 2, 5 et 9. Ce parcours
          couvre le principe de contraction de Banach et le th\'eor\`eme de Brouwer
          avec leurs principales applications.
    \item \textbf{Parcours avanc\'e} (M2)~: Ajouter les Chapitres 3, 6, 7 et 8.
          Ce parcours compl\`ete la formation par les extensions m\'etriques, le
          th\'eor\`eme de Schauder, le th\'eor\`eme de Kakutani et les espaces ordonn\'es.
    \item \textbf{Parcours recherche} (Doctorat)~: L'int\'egralit\'e du cours,
          avec une attention particuli\`ere aux Chapitres 4, 10 et 11 qui touchent
          \`a des th\`emes de recherche actuels.
\end{itemize}

\section*{Remerciements}

Ce cours doit beaucoup aux ouvrages classiques de la discipline, notamment ceux
de Granas et Dugundji \cite{granas2003}, Goebel et Kirk \cite{goebel1990},
Zeidler \cite{zeidler1986}, Agarwal, Meehan et O'Regan \cite{agarwal2001},
et Smart \cite{smart1974}. Nous remercions \'egalement les nombreux coll\`egues
et \'etudiants dont les remarques ont contribu\'e \`a am\'eliorer ce texte.

\section*{Comment utiliser ce cours}

Chaque chapitre contient~:
\begin{itemize}
    \item Des \textbf{d\'efinitions} num\'erot\'ees, pr\'ecises et motiv\'ees;
    \item Des \textbf{th\'eor\`emes} avec d\'emonstrations compl\`etes;
    \item Des \textbf{exemples} illustratifs et des \textbf{contre-exemples}
          montrant la n\'ecessit\'e des hypoth\`eses;
    \item Des \textbf{remarques} \'eclairant les points subtils;
    \item Des \textbf{exercices} de difficult\'e vari\'ee, certains avec indications.
\end{itemize}

Les encadr\'es sp\'eciaux (formules cl\'es, attention, intuition, algorithme)
permettent de rep\'erer rapidement les informations essentielles.

\bigskip
\hfill \textit{L'auteur}

\hfill \textit{Mars 2026}
